% SHARD-ID: RC-04O-GRADED-TOYS-EXITS
% SHARD-TITLE: The Phantasm VII: a Graded Tensor Network for the Cusp Toys II. The Even Exit and the Funnel, the Glued Toy with Odd Letters, the Frobenius Channel of the Arithmetic Metric, and the Level Tower
% SHARD-SUMMARY: On a regular diagram with cusps and funnels the constant function is bound on the cusps and outgoing on the funnels at the same z = q^{-1/2}, so it is neither a resonance nor a bound state; a funnel bolted onto an arithmetic core moves the Hasse--Weil quartet off the Weil circle at first order; and on any graded bond a channel with only even letters has two stationary densities, so a unique mixed stationary state with the pole even and the zeros odd requires an odd (fermionic) letter, realised by a glued network whose sectors exit through different blocks and are protected exactly because they are energy-orthogonal.
% SHARD-SUMMARY: The arithmetic-metric completion with one exit per zero mode is the channel r^2 Ad(U) + (1 - r^2) Omega Tr with U^2 the normalised Frobenius, a one-gap quantum expander whose characteristic function is diagonal in Blaschke factors; no self-adjoint diagram has it, because each constant exit direction of a Hermitian core sees a conjugation-closed set of resonances.
% SHARD-SUMMARY: Gamma_0(N) with trivial character has only oldform Eisenstein series and sees only zeta; the Dirichlet L-zeros live on Gamma_1(N), where, under H-ARITH-1, the graded bond is one even vacuum line from the trivial character and odd-only L-zero blocks from the others, graded by the diamond operators.
% SHARD-KEYWORDS: funnel, constant mode, half-outgoing, even exit, odd letter, fermionic letter, glued toy, energy-orthogonality, shared exit, persistence, arithmetic metric, Frobenius channel, quantum expander, Blaschke, characteristic function, Hermitian core, Klein group, Gamma_0(N), Gamma_1(N), oldforms, Dirichlet L-function, diamond operator, H-ARITH-1
% SHARD-NOTES: none
% SHARD-SCRIPTS: none

\section{The Phantasm VII: a Graded Tensor Network for the Cusp Toys II. The Even Exit, the Glued Toy, the Frobenius Channel, the Level Tower}
\label{sec:graded-toys-exits}

Same lanes and files as \cref{sec:graded-toys-network}. This shard answers items (i), (iv) and (iii) of
\cref{obs:elliptic-cavity-next}.

\subsection{The even exit: what a funnel does}

\begin{proposition}[The constant function on a diagram with cusps and funnels is half-outgoing]
\label{prop:constant-mode-mixed-sheets}
\claimstatus{prop:constant-mode-mixed-sheets}{sketched}
Let a diagram with core $T_X$, cusp rays and funnel rays ($S(f_{k+1}) = S(f_k)/\qs$, $S(f_kf_{k+1}) = S(f_{k+1})$)
be $(\qs+1)$-regular at every vertex, and $C_c$, $C_f$ the cusp and funnel junction weights
(\cref{prop:funnel-self-energy}). The restriction $x$ of the constant function to the core, with all entries
nonzero, satisfies
\[
\big[(1+z^2)I - zT_X - z^2C_c - \qs^{-1}C_f\big]x = 0\qquad\text{at } z = \qs^{-1/2},
\]
bound sheet on every cusp, outgoing sheet on every funnel; it is a null vector of the outgoing matrix of
$p_{\mathrm{funnel}}$ iff $C_c = 0$ and of the bound matrix iff $C_f = 0$. For the one-vertex regular diagram with
cusp weight $a$ and funnel weight $\qs+1-a$, $p = z^2-(1+a(\qs-1))/\qs$: for $0\le a<1$ the Perron root is a resonance
at $z^2 = (1+a(\qs-1))/\qs\in[1/\qs,1)$, at $\qs^{-1/2}$ iff $a = 0$; $a = 1$ is a threshold; for $1<a\le\qs+1$ it is a
visible bound state with $\vartheta^2 = \qs/(1+a(\qs-1))$, at $\qs^{-1/2}$ iff $a = \qs+1$. A light cusp leaves the
Perron-type mode outgoing, but at a radius other than $\qs^{-1/2}$.
\end{proposition}

\begin{lemma}[First-order motion of the resonances under an added funnel]
\label{lem:funnel-first-order}
\claimstatus{lem:funnel-first-order}{sketched}
Adding a funnel of weight $f$ at vertex $v$ gives exactly $p_f = p-(f/\qs)\,p^{(v)}$, $p^{(v)}$ the resonance
polynomial of the core with $v$ deleted, so a simple root moves by $dz_i/df = \qs^{-1}p^{(v)}(z_i)/p'(z_i)$ and
$d|z_i|/df = \re(\bar z_i\,dz_i/df)/|z_i|$. Since $P_v$ is real and diagonal, the roots of $p_f$ stay closed under
$z\mapsto-z$ and $z\mapsto\bar z$ for every $f$, so the Hasse--Weil quartet remains one Klein orbit
$\{z,-z,\bar z,-\bar z\}$ and hence equimodular as long as it has four distinct non-real, non-imaginary elements: an
added funnel moves the common radius off $\qs^{-1/4}$ but cannot break the equal-modulus property until the orbit
splits. At genus one, $\mathrm{RH}(Y)$ on $K_{\mathrm{HW}}$ is therefore a consequence of the bipartite sign and reality
whenever the quartet is a single orbit; Weil's content for D2 is the radius, pinned by the product identity of
\cref{prop:nonzero-root-product}. At genus two, two orbits can have different moduli.
\end{lemma}

\begin{numerical}[A funnel bolted onto D2]
\label{num:funnel-on-d2}
\claimstatus{num:funnel-on-d2}{numerical}
\texttt{scripts/graded\_toys.py}, ledger D32--D37: the constant-mode identity of \cref{prop:constant-mode-mixed-sheets}
on three regular cusp-plus-funnel diagrams, both ``iff'' clauses; the one-vertex formula at $\qs = 2,3,5$; on D2 with a
funnel of weight $f$ at each core vertex, $d|z_i|/df$ at $f = 0$ equal for all four Hasse--Weil roots at every vertex
($+0.0630672$ at $A$, $+0.1261345$ at $B$, $-0.0210224$ at $C$, $-0.1681793$ at $D$, $-0.0630672$ at $E$ and $F$), the
four moduli equal to $10^{-15}$ at $f = 1/5$ (radii $0.853128$, $0.864176$, $0.836761$, $0.807041$, $0.827557$), the
orbit splitting into two moduli at $f = 1$ ($D$: $0.618034$, $0.707107$), $f = 2$ ($E$, $F$), $f = 5$ ($A$), not up to
$f = 5$ at $B$, $C$; the Perron parameter decreasing strictly from $\qs^{-1/2}$ ($0.706576$ at $f = 0.01$, $0.138712$
at $f = 100$, funnel at $B$), the Perron pair a visible bound state for every $f$ tried; with the cusp of D2 replaced
by a funnel, the root motions.
\end{numerical}

So item (i) is negative on one connected regular diagram: a cusp of weight $a>1$ keeps the constant mode bound, a
light cusp makes a Perron-type mode outgoing at the wrong radius, and a funnel on the arithmetic core costs the
Weil radius, though not, at genus one, the equal-modulus property. The graded network says what is needed instead.

\subsection{The glued toy: odd letters are necessary and sufficient}

\begin{lemma}[Even letters give two stationary densities]
\label{lem:even-letters-two-fixed-points}
\claimstatus{lem:even-letters-two-fixed-points}{sketched}
Let $H = H_+\oplus H_-$ with both summands nonzero and $\rdens\mapsto\sum_iA_i\rdens A_i^{*}$ trace preserving with
every $A_i$ even. Then there are stationary densities supported on $H_+$ and on $H_-$. Hence a parity-preserving
channel with a unique stationary density has, in every Kraus representation by homogeneous letters, an odd letter.
\end{lemma}

Even letters and their adjoints preserve each sector, so the channel restricts to a trace-preserving map on
each $B(H_\pm)$, which has a fixed density (\texttt{proofs.md}, G4(d)(iii)). In the fermionic-MPS language of
\cref{def:cmps-ring-vector} an odd letter is a fermionic physical species: the parity-crossing rebound is a
fermionic jump, the ``parity jump'' of the ideation ledger.

\begin{theorem}[The glued toy]
\label{thm:glued-graded-toy}
\claimstatus{thm:glued-graded-toy}{sketched}
Let $(Z_\pm,J_\pm)$ be $h_\pm$-exit contractions on $K_\pm$, $K = K_+\oplus K_-$ with $P = 1\oplus(-1)$,
$Z = Z_+\oplus Z_-$, $J = J_+\oplus J_-$ into $\mathbb C^{h_+}\oplus\mathbb C^{h_-}$, and $\reb = \reb_+\oplus\reb_-$
an even density with $\Tr\reb_\pm>0$. (i) The letters $A_0 = Z$, $A_{(a,k)} = \sqrt{p_k}|\omega_k\rangle\langle j_a|$
are homogeneous of parity $\epsilon(\omega_k)\epsilon(j_a)$, the cross letters odd; the channel is a
trace-preserving graded transfer channel. (ii) It has a unique stationary density
$\rdens_\infty = \bar t^{\,-1}\sum_mZ^m\reb Z^{*m}$, even, of rank at least two, with $\rdens_\infty|_{K_\pm}\ge\bar t^{\,-1}\reb_\pm$;
it attracts every density iff the holding law is aperiodic. (iii) On the odd sector the channel is
$X\mapsto ZXZ^{*}$ (the exit matrix of an odd $X$ is block-off-diagonal, of trace zero): the odd coherences
$|k_n\rangle\langle k_m|$, $k_n\in K_-$, $k_m\in K_+$, are eigen-operators with eigenvalues $z_n\bar z_m$ for every
fixed reset, of modulus $r_+r_-$ under $\mathrm{RH}(Y_+)$, $\mathrm{RH}(Y_-)$: rates add. (iv) The odd part of the
twisted ring norm is $\str\Edb^L-\Tr\Esec{0}^L = -2\re[(\Tr Z_-^L)\,\overline{\Tr Z_+^L}]$. (v) Instance: $K_+$ the
model of APW's tree, the one-vertex core with one funnel of weight $\qs+1$ ($c_+^2 = (\qs+1)/\qs$,
$p_+ = z^2-1/\qs$, no bound states, $\dim K_+ = 2$, $\operatorname{spec}Z_+ = \{\pm\qs^{-1/2}\}$), $K_-$ the
six-dimensional model of D2: odd eigenvalues $\pm\qs^{-1/2}\bar z_n$ and conjugates, of modulus $\qs^{-3/4}$ on
the Hasse--Weil part and $0$ on the delay part, and $\Tr Z_+^L = \qs^{-L/2}(1+(-1)^L)$.
\end{theorem}

\begin{proposition}[Shared exit: protection is energy-orthogonality]
\label{prop:shared-exit-tradeoff}
\claimstatus{prop:shared-exit-tradeoff}{sketched}
For a fixed reset and eigenvectors $k_n,k_m$ of $Z$ with simple eigenvalues, $X = |k_n\rangle\langle k_m|$ has
exit trace $\Tr(JXJ^{*}) = \langle a_m,a_n\rangle = (1-\bar z_mz_n)\langle k_m,k_n\rangle$. If it vanishes, $X$ is an
eigen-operator with eigenvalue $z_n\bar z_m$; if not, $X$ is not an eigen-operator (residue
$\langle a_m,a_n\rangle\reb$), and the copy of the eigenvalue it carried persists iff $\langle l_n|\reb|l_m\rangle = 0$ for
the dual eigenvectors (\cref{prop:rebound-persistence-secular}); on a real core the value is degenerate in $\Ad(Z)$
and keeps its other copies. On a
connected core with a shared exit the modes are not orthogonal, so the zeros are not protected once the even
mode leaks; the glued toy is protected because its sectors exit through different blocks.
\end{proposition}

Conclusion for item (i): the object with pole even, zeros odd and a unique mixed stationary state is not a
graph of groups (\cref{cit:apw-cusps-pic-no-funnels}) but a graded bond gluing an even funnel diagram to an odd
cusp diagram through a parity-crossing rebound: the ``hedgehog'' has an even spike and odd spikes that meet
only in the rebound, and the fermionic letter that joins them is exactly what makes the stationary state unique.

\subsection{The Frobenius channel: item (iv)}

\begin{theorem}[The arithmetic-metric renewal channel is a one-gap quantum expander with diagonal characteristic function]
\label{thm:frobenius-channel-expander}
\claimstatus{thm:frobenius-channel-expander}{sketched-conditional}
Let $U$ be unitary on a finite-dimensional $K$, $0<r<1$, $Z' = rU$, $J' = \sqrt{1-r^2}\,1_K$ (one exit per mode),
and $\reb$ a density. The renewal channel $\mathcal E'(\rdens) = r^2U\rdens U^{*}+(1-r^2)\Tr(\rdens)\reb$ is CPTP with
unique stationary density $\rdens_\infty = (1-r^2)\sum_mr^{2m}U^m\reb U^{*m}$, equal to $\reb$ iff $[U,\reb] = 0$; on
traceless operators $\mathcal E' = r^2\Ad(U)$, so every relaxation eigenvalue has modulus exactly $r^2$. The
characteristic function of $Z'$ is $\Theta'(z) = (1-zrU^{*})^{-1}(z-rU) = \operatorname{diag}\big((z-z_a)/(1-\bar z_az)\big)$
in the eigenbasis of $U$: the $a$-th exit sees the single mode $z_a$. For $K_{\mathrm{HW}}$ of D2 in the symmetric
arithmetic metric (\cref{thm:arithmetic-metric}), $r = \qs^{-1/4}$ and $U^2\simeq(\Frob/\sqrt{\qs})\otimes1_2$
as multisets; the relaxation moduli are all $\qs^{-1/2}$.
\end{theorem}

\begin{proposition}[No self-adjoint diagram has one-mode exits; the Klein group of the odd sector]
\label{prop:no-selfadjoint-one-mode-exits}
\claimstatus{prop:no-selfadjoint-one-mode-exits}{sketched}
For a Hermitian core with standard rays, $\mathbf S(\bar z) = \mathbf S(z)^{*}$, hence $e^{*}\mathbf S(\bar z)e =
\overline{e^{*}\mathbf S(z)e}$ for every constant $e\in\mathbb C^h$; if $\mathbf S$ is diagonalised by a constant
unitary, each diagonal channel has a conjugation-closed zero set. So the diagonal $\Theta'$ of
\cref{thm:frobenius-channel-expander} with the four non-real Hasse--Weil modes in four entries is not the
scattering matrix of any self-adjoint diagram; the finest constant-basis splitting of the quartet is into two
conjugation-closed pairs. On $K_{\mathrm{HW}}$ the bipartite sign and complex conjugation generate a Klein
four-group acting simply transitively on the four modes; the symmetric arithmetic metric is the unique diagonal
metric it fixes, and in it the four one-mode exits are permuted by the group.
\end{proposition}

Answer to item (iv): the object with four cusps seeing the four modes separately is the Frobenius channel,
the odd disc part of $E_S$ (\cref{thm:scattering-superdeterminant}) completed by a fixed reset; it is a
Lindblad-level object, not the wave equation of a graph, and its uniform gap is the metric's, not a theorem
about D2.

\subsection{The level tower: item (iii)}

\begin{observation}[$\Gamma_0(N)$ with trivial character sees only $\zeta$; the $L$-zeros live on $\Gamma_1(N)$]
\label{obs:level-tower-correction}
\claimstatus{obs:level-tower-correction}{sketched}
For $\Gamma_0(N)<\SL(\mathbb Z)$ with trivial nebentypus and $N$ squarefree, the $2^{\omega(N)}$ functions
$E(dz,s)$, $d\mid N$, are $\Gamma_0(N)$-invariant and span the Eisenstein space (there are no Eisenstein newforms
with trivial character at squarefree level), so every scattering-matrix entry is a rational function of the
$\pp^{-s}$, $\pp\mid N$, times $\cxi(2s-1)/\cxi(2s)$, and the scattering determinant is
$(\cxi(2s-1)/\cxi(2s))^{2^{\omega(N)}}$ times a rational function: no Dirichlet $L$-function. The $L$-ratios
$L(2s-1,\chi)/L(2s,\chi)$ enter for $\Gamma_1(N)$, equivalently $\Gamma_0(N)$ with a nontrivial nebentypus
\cite{huxley1984scattering,hejhal1983selberg2}; \cref{cit:kk-level-constant-term} is the trivial-character
instance over $\Fq[T]$ at one cusp. Item (iii) and \cref{conj:galois-graded-bond} should name $\Gamma_1(N)$.
Not byte-cited; the sources are not in the repository.
\end{observation}

\begin{assumption}[H-ARITH-1]
\label{asm:h-arith-gamma1}
\claimstatus{asm:h-arith-gamma1}{assumed}
For $\Gamma_1(N)<\mathrm{GL}_2(\Fq[T])$ the scattering matrix is diagonalised by the Dirichlet characters $\chi$ of
$(\Fq[T]/N)^\times$, the $\chi$-channel being a monomial in $z$ times $L(2s-1,\chi)/L(2s,\chi)$ with $L$ the
completed $L$-function (Euler factor at $\infty$ included), and the diamond operators act on the $\chi$-channel by
$\chi$.
\end{assumption}

\begin{proposition}[The graded bond of a character channel is odd-only]
\label{prop:character-channel-graded-bond}
\claimstatus{prop:character-channel-graded-bond}{sketched-conditional}
For a nontrivial \emph{primitive} $\chi$ mod $N$ over $\Fq[T]$, $L(\uz,\chi) = \sum_{f\ \mathrm{monic}}\chi(f)\uz^{\deg f}$ is
a polynomial of degree $\deg N-1$, equal to $\det(1-\uz\Frob_\chi)$ for odd $\chi$ and to $(1-\uz)\det(1-\uz\Frob_\chi)$
for even $\chi$ (the factor of the place at $\infty$), with every eigenvalue of $\Frob_\chi$ of modulus
$\sqrt{\qs}$ \cite{rosen2002function}. With $E_\chi = (\Frob_\chi)_-\oplus\Pi\Frob_\chi(-1)$ and $w = \qs^{-2s}$, the
completed ratio is $L(2s-1,\chi)/L(2s,\chi) = \sdet(1-wE_\chi)$, whose disc part is $\netmult(\alphaw{\chi,i}) = -1$
at the eigenvalues of $\Frob_\chi$ and has no even entry; the trivial character carries
$\zeta_{\Fq(T)}(2s-1)/\zeta_{\Fq(T)}(2s) = \sdet(1-wE_S)$ at $g = 0$, disc part $\netmult(1) = +1$. Under
\cref{asm:h-arith-gamma1} the disc part of the level-$N$ graded bond is
$(1)_+\oplus\bigoplus_{\chi\ \mathrm{primitive},\ \chi\ne1}(\alphaw{\chi,1},\dots,\alphaw{\chi,d_\chi})_-$ over the conductors
dividing $N$, the trivial zero of each even $\chi$ split off as the vacuum line is split off the trivial channel, the
diamond operator $\langle a\rangle$ acting on the $\chi$-block by $\chi(a)$: over $\Fq(T)$ every zero lives in the level
tower, and the vacuum in the trivial channel alone. Imprimitive characters add the Euler factors of the primes
dividing $N$ but not the conductor, with roots of modulus one; these, and the trivial zeros if the incomplete $L$
appears, are odd modes at the Perron radius $|z| = \qs^{-1/2}$.
\end{proposition}

\begin{numerical}[{Character channels over $\mathbb F_3[T]$}]
\label{num:character-channels}
\claimstatus{num:character-channels}{numerical}
\texttt{scripts/graded\_toys.py}, ledger D46--D51: all $25$ nontrivial characters mod $T^3-T-1$ (primitive, $\deg L = 2$,
$12$ even with the factor $1-\uz$) and all $15$ mod $T(T^2+1)$ ($7$ primitive); every nontrivial root of every primitive
character has $|\alphaw{}| = \sqrt3$; the Euler product to order $\uz^5$; $\det(1-\uz\Frob_\chi) = L(\uz,\chi)$; the
superdeterminant identity as rational functions for all $40$ characters; no even disc eigenvalue in any channel; the
imprimitive characters' modulus-one roots; the trivial channel at $g = 0$. Also this shard's glued toy (D38--D45:
parities of the letters, uniqueness by the eigenvalue count at $1$, both sectors populated, the odd eigenvalues
$\pm\qs^{-1/2}\bar z_n$, the twisted ring-norm identity, the segment for a sector-preserving reset, the exact vanishing
$\langle a_m,a_n\rangle = 0$ across blocks, the shared-exit residue $\langle a_m,a_n\rangle\reb$ on a random three-vertex
core) and the Frobenius channel (G5(a,b) in full; conjugation-closed disc zero sets for real and complex exit
directions on D2 with a second cusp at each vertex).
\end{numerical}

\begin{observation}[What the next round should do]
\label{obs:graded-toys-next}
\claimstatus{obs:graded-toys-next}{sketched}
(i) Compute the $\Gamma_1(N)$ (or nebentypus) Eisenstein constant terms over $\Fq[T]$ for one $N$ of degree
three and settle \cref{asm:h-arith-gamma1}, including the completed-versus-incomplete question; the quotient
graph $\Gamma_1(N)\backslash\mathcal T$ (Gekeler--Nonnengardt) is the diagram. (ii) The glued toy as a
fermionic MPS proper: write the odd letters as fermionic species of \cref{def:cmps-ring-vector} and compute
the antiperiodic and periodic closures with the sign-twisted transfer generator. (iii) Ask whether the Frobenius
channel of \cref{thm:frobenius-channel-expander} is the transfer channel of a Weil-representation or graded
Harrow expander (\cref{sec:graded-ramanujan}) at $\pp = 2$: its letters are $rU$ and the reset, its relaxation
spectrum is $r^2\operatorname{spec}\Ad(U)$. (iv) A non-self-adjoint core (gain and loss at the junctions)
realising the diagonal $\Theta'$, as the Lindblad-level ``hedgehog'' the arithmetic metric asks for.
\end{observation}
